\chapter{Общая бимодульная кривизна и относительный Грам-вес}
\label{ch:v8-bimodule-common-curvature-relative-weight-gate}

\section{Честный общий след}

После бимодульного замыкания поле переноса имеет размерность $20$ над
$\mathbb C$. Его завершённый вещественной структурой рёберный момент представляется самосопряжённой
кривизной
\begin{equation}
 \Omega_E=\operatorname{diag}(m_1,\ldots,m_{20},-m_1,\ldots,-m_{20}),
 \qquad \frac12\Tr\Omega_E^2=S_E.
 \label{eq:v8-common-edge-curvature}
\end{equation}
Физическая Грам-кривизна живёт на двух концевых углах:
\begin{equation}
 \Omega_G(A)=\operatorname{diag}
 \bigl(AA^*-A_0A_0^*,\;A^*A-A_0^*A_0\bigr),
 \qquad \frac12\Tr\Omega_G^2=S_G.
 \label{eq:v8-common-gram-curvature}
\end{equation}
Поэтому незавешенная общая кривизна
\begin{equation}
 \Omega=\Omega_E\oplus\Omega_G
 \label{eq:v8-unweighted-common-curvature}
\end{equation}
даёт строго $S_E+S_G$, то есть $\beta=1$. удвоение вещественной структуры и физический
полуслед умножают оба слагаемых одинаково и отношения не меняют.

Чтобы получить проходящий представитель $\beta=1/2$, необходимо заменить
$\Omega_G$ на $\Omega_G/\sqrt2$. Такая замена эквивалентна явной вставке
искомой относительной метрики и потому не является её выводом.

\section{Проверка скрытого индекс-один коннектора}

Размерности $20$ и $21$ создают естественное подозрение, что после
бимодульного расширения мог появиться почти унитарный коннектор
\begin{equation}
 J:E_t\oplus E_s\longrightarrow\mathcal T_{\rm bimod},
 \qquad \dim(E_t\oplus E_s)=21.
 \label{eq:v8-index-one-connector-candidate}
\end{equation}
Для физической калибровочной ковариантности он обязан удовлетворять
\begin{equation}
 R_{\mathcal T}(X)J=J R_E(X)
 \label{eq:v8-connector-intertwiner-equation}
\end{equation}
для всех двенадцати генераторов.

Полная линейная система имеет $13$ комплексных решений, но их общий
кообраз конечного отображения имеет размерность только $9$. Следовательно,
\begin{equation}
 \rank J\leq9<20
 \label{eq:v8-connector-rank-bound}
\end{equation}
для каждого ковариантного интертвинера. Почти унитарного коннектора ранга
$20$ нет. Полярная коизометрия $A_0:\mathbb C^{11}\to\mathbb C^{10}$
связывает два концевого угла, но не превращается в отображение из
$21$ концевой компоненты в $20$ коэффициентов поля переноса.

\section{Центральная свобода и гессиан}

На общем носителе размерности $40+21=61$ проекторы $P_E$ и $P_G$ остаются
центральными относительно блочно-диагональной кривизны. Поэтому всё
семейство
\begin{equation}
 G_\eta=P_E+\eta P_G,
 \qquad \eta>0,
 \label{eq:v8-common-central-metric-family}
\end{equation}
допустимо. Незавешенный выбор $\eta=1$ не привилегирован внутренней
симметрией и, более того, даёт провальную исходную сигнатуру
$(38,0,2)$. При $\beta=1/2$ сохраняется требуемая сигнатура
$(20,0,20)$; при $\beta=2/3$ возникают двенадцать нулевых мод.

\begin{theorem}[Запрет незавешенной общей кривизны]
На бимодульно замкнутом модуле переноса одна незавешенная прямая сумма
рёберно-ходжева и грамова метрика концевых секторов кривизн фиксирует $\beta=1$, а не проходящее
значение. вещественный полуслед этого отношения не меняет. Ковариантный линейный
коннектор $21\to20$ полного ранга отсутствует: максимальный возможный ранг
равен $9$.
\end{theorem}

\begin{nogocriterion}[Запрет скрытой нормировки]
Нельзя объявлять $\beta=1/2$ выведенным посредством записи
$\Omega_E\oplus\Omega_G/\sqrt2$: множитель $1/\sqrt2$ является ровно тем
свободным относительным весом, происхождение которого проверяется. Нельзя
также использовать одну лишь близость размерностей $20$ и $21$ как
основание полярного подъёма: калибровочный интертвинер ранга $20$ не существует.
\end{nogocriterion}

\begin{protocol}
\begin{enumerate}
 \item общий кривизностный носитель: \textbf{$40+21=61$};
 \item незавешенный общий след: \textbf{$\beta=1$};
 \item вещественный полуслед меняет отношение: \textbf{нет};
 \item проходящий представитель: \textbf{$\beta=1/2$};
 \item его происхождение: \textbf{не выведено};
 \item пространство ковариантных интертвинеров: \textbf{$13$ комплексных измерений};
 \item максимальный ранг интертвинера $21\to20$: \textbf{$9$};
 \item полноранговый коннектор: \textbf{отсутствует};
 \item центральное семейство $G_\eta$: \textbf{сохраняется};
 \item сигнатура при $\beta=1$: \textbf{$(38,0,2)$};
 \item полный родительское действие: \textbf{не получен};
 \item следующий гейт: \textbf{нелинейный граница инцидентности коннектор}.
\end{enumerate}
\end{protocol}

Машинный сертификат сохранён в
\path{s2t_v8_bimodule_common_curvature_relative_weight_gate_results.json}.